\SourceLeaf{335}%
\Article{34}
Von den sieben Coëfficienten (1), (2) u.\,s.\,w. ist der letzte constant, nemlich
\[
\log(7)=7{,}6287228032(-20)
\]
und
\[
\log3(7)=8{,}1058440580(-20)
\]
die übrigen werden, sobald bestimmte Werthe für die Dimensionen des Ellipsoids gewählt sind, Functionen der Breite \(B\), und lassen sich also in eine Tafel bringen, deren Argument \(B\) ist. Steht eine solche Tafel zu Gebote, so ist die Rechnung nach dieser Methode für das Ellipsoid eben so bequem, wie die Rechnung für die Kugel.

Ich füge am Schlusse dieser Abhandlung eine solche Tafel für die Zone von \(51^0\) bis \(54^0\) bei, in welcher die Werthe von \(B\) von Minute zu Minute fortschreiten, und bemerke dazu folgendes.

Von den Ellipsoidelementen ist die Tafel nur in so fern abhängig, als darin eine bestimmte Abplattung oder ein bestimmter Werth von \(e\) zum Grunde gelegt ist, derjenige nemlich, welchen die letzte von \textsc{Bessel} ausgeführte Rechnung ergeben hat, und der auch der der ersten Abhandlung beigefügten Tafel zum Grunde liegt (s. Art. 5). Damit der Zahlenwerth von \(a\) bloss von der Abplattung abhängig werde, ist als Einheit nicht die Toise oder ein sonstiges willkürliches Maass angenommen, sondern der zehnmillionste Theil des Erdmeridians, wonach also \(a\) unmittelbar durch \(e\) vermittelst der Gleichung
\[
\begin{aligned}
\pi a(1-\tfrac14ee-\frac{1\mathord{.}3}{4\mathord{.}16}e^4-\frac{1\mathord{.}3\mathord{.}15}{4\mathord{.}16\mathord{.}36}e^6\\
{}-\frac{1\mathord{.}3\mathord{.}15\mathord{.}35}{4\mathord{.}16\mathord{.}36\mathord{.}64}e^8
-\frac{1\mathord{.}3\mathord{.}15\mathord{.}35\mathord{.}63}{4\mathord{.}16\mathord{.}36\mathord{.}64\mathord{.}100}e^{10}-\text{u.\,s.\,w.})\\
=20000000
\end{aligned}
\]
deren Gesetz offenbar ist, gefunden werden kann, oder vermittelst der ihr gleichgeltenden
\[
a=\frac{20000000}{\pi}(1+\tfrac14ee+\tfrac7{64}e^4+\tfrac{15}{256}e^6+\tfrac{579}{16384}\mathord{.}e^8+\tfrac{1515}{65536}e^{10}+\text{u.\,s.\,w.})
\]
Man findet so, mit jenem Werthe von \(e\).
\[
\begin{aligned}
a&=6376851{,}447\\
\log a&=6{,}8046062999
\end{aligned}
\]
Es versteht sich, dass bei Anwendung unsrer Tafel auch \(r\) erst in derselben Einheit ausgedrückt sein muss; um dies zu erreichen, wird man (gemäss dem von
